\section{Syntax-Directed Definition}
\label{sec:ch14-sdd}

\subsection{Definisi}

\begin{definisi}[SDD]
Syntax-directed definition (SDD) adalah tata bahasa dengan atribut
dan aturan semantik. Setiap produksi \(A \to \alpha\) memiliki aturan
untuk menghitung atribut simbol-simbolnya.
\end{definisi}

SDD mendeskripsikan \emph{apa} yang dihitung, bukan \emph{kapan}
perhitungan dilakukan. Urutan evaluasi ditentukan oleh
ketergantungan antar-atribut.

\subsection{Contoh SDD untuk Ekspresi}

\begin{table}[H]
\centering
\caption{SDD sederhana untuk tipe ekspresi}
\label{tab:ch14-sdd}
\begin{tabularx}{\textwidth}{@{}l X@{}}
\toprule
\textbf{Produksi} & \textbf{Aturan semantik} \\
\midrule
\(E \to E_1 + E_2\) &
\(E.type = \text{arith}(E_1.type, E_2.type)\) \\
\(E \to E_1 * E_2\) &
\(E.type = \text{arith}(E_1.type, E_2.type)\) \\
\(E \to \code{id}\) &
\(E.type = \text{lookup}(\code{id}.name)\) \\
\(E \to \code{INT\_LIT}\) & \(E.type = \text{int}\) \\
\(E \to \code{FLOAT\_LIT}\) & \(E.type = \text{float}\) \\
\bottomrule
\end{tabularx}
\end{table}

Aturan \code{arith} mengembalikan \code{float} jika salah satu
operand \code{float}, \code{int} jika keduanya \code{int} atau
\code{char}, dan \code{error} jika operand tidak numerik.

\subsection{Aturan untuk Pernyataan}

Ekspresi saja tidak cukup. Sub-CPMK 1.5 meminta aturan untuk ekspresi
dan pernyataan. Tabel \tabref{tab:ch14-sdd-stmt} melengkapi tabel
sebelumnya dengan aturan untuk assignment, if, while, blok, dan
deklarasi.

\begin{table}[H]
\centering
\caption{Tabel aturan semantik untuk pernyataan}
\label{tab:ch14-sdd-stmt}
\begin{tabularx}{\textwidth}{@{}l X X@{}}
\toprule
\textbf{Produksi} & \textbf{Aturan semantik} &
\textbf{Keluaran yang diuji} \\
\midrule
\(S \to \code{id} = E\) &
\(S.type = \text{check\_assign}(\text{lookup}(\code{id}),\allowbreak\
E.type)\) &
tipe \(S\) atau \code{error} \\
\(S \to \code{if}\ (E)\ S_1\) &
\(\text{check\_bool}(E.type);\ S.type = S_1.type\) &
tipe \(S\) \\
\(S \to \code{while}\ (E)\ S_1\) &
\(\text{check\_bool}(E.type);\ S.type = S_1.type\) &
tipe \(S\) \\
\(S \to \{\,S_1\ S_2\,\}\) &
\(S.type = \text{seq}(S_1.type, S_2.type)\) &
tipe \(S\) \\
Deklarasi \(D\) &
\(\text{declare}(\code{id}, \code{type})\) &
entri tabel simbol \\
\bottomrule
\end{tabularx}
\end{table}

\subsection{Menguji pada Pohon Sintaks}

Aturan pada tabel diuji dengan menelusuri pohon sintaks. Untuk program
\begin{lstlisting}[language=C]
int a = 2;
float b = 3.5;
a = b + 1;
\end{lstlisting}
hasil penelusuran setiap simpul dicatat pada
\tabref{tab:ch14-uji}.

\begin{table}[H]
\centering
\caption{Hasil pengujian aturan semantik pada pohon sintaks}
\label{tab:ch14-uji}
\begin{tabularx}{\textwidth}{@{}l X X@{}}
\toprule
\textbf{Simpul} & \textbf{Aturan} & \textbf{Hasil} \\
\midrule
Deklarasi \(a\) & \(\text{declare}(a, \code{int})\) &
entri \(a:\code{int}\) \\
Deklarasi \(b\) & \(\text{declare}(b, \code{float})\) &
entri \(b:\code{float}\) \\
Ekspresi \(b + 1\) & \(\text{arith}(\code{float}, \code{int})\) &
\code{float} \\
Assignment \(a = b + 1\) &
\(\text{check\_assign}(\code{int}, \code{float})\) &
\code{error}: perlu cast \\
\bottomrule
\end{tabularx}
\end{table}

Tabel ini menunjukkan bahwa aturan dapat diuji sebelum
diimplementasikan. Setiap baris harus merujuk simpul pohon sintaks
yang konkret, bukan hanya nama aturan.

\subsection{S-Attributed dan L-Attributed}

Atribut synthesized dihitung dari anak, sedangkan atribut inherited
dihitung dari orang tua atau saudara; rincian dan contohnya ada pada
Bagian 14.4. SDD S-attributed hanya memakai atribut synthesized,
sehingga dapat dievaluasi dengan traversal post-order. SDD
L-attributed memakai atribut inherited yang hanya bergantung pada
simbol di kiri atau orang tua, sehingga dapat dievaluasi dengan
traversal kiri-ke-kanan.

\begin{table}[H]
\centering
\caption{Perbandingan S-attributed dan L-attributed}
\label{tab:ch14-s-l}
\begin{tabularx}{\textwidth}{@{}lXX@{}}
\toprule
\textbf{Aspek} & \textbf{S-attributed} & \textbf{L-attributed} \\
\midrule
Atribut & Synthesized & Synthesized dan inherited \\
Evaluasi & Post-order & Kiri-ke-kanan \\
Implementasi & Mudah & Perlu konteks \\
Contoh & Tipe ekspresi & Lingkungan scope \\
\bottomrule
\end{tabularx}
\end{table}

\subsection{Evaluasi}

Evaluasi SDD dapat dilakukan:
\begin{itemize}[leftmargin=*]
  \item saat parsing, jika aturan sudah cukup;
  \item setelah AST terbentuk, dengan traversal;
  \item dengan alat khusus yang membangun dependency graph.
\end{itemize}

Pada proyek Subset-C, evaluasi dilakukan setelah AST terbentuk.
Pendekatan ini memisahkan parsing dari analisis dan memudahkan
pengujian.
